\documentclass[12pt]{article}
\usepackage{graphicx} % Required for inserting images
\usepackage[margin=0.5in]{geometry}
\usepackage{amsmath, amssymb}
\usepackage{mathrsfs}


\usepackage{soul}


% Dr. David Brown Commands
\newcommand{\ds}{\displaystyle}
\newcommand{\R}{\mathbb{R}}
\newcommand{\N}{\mathbb{N}}
\newcommand{\Q}{\mathbb{Q}}
\newcommand{\C}{\mathbb{C}}


% Commands to get script r
\usepackage{calligra}
\DeclareMathAlphabet{\mathcalligra}{T1}{calligra}{m}{n}
\DeclareFontShape{T1}{calligra}{m}{n}{<->s*[2.2]callig15}{}
\newcommand{\scripty}[1]{\ensuremath{\mathcalligra{#1}}}

% Unit commands
\newcommand{\degree}{\text{°}}
\newcommand{\unitSpeed}{\frac{\text{m}}{\text{s}}}
\newcommand{\unitAcceleration}{\frac{\text{m}}{\text{s}^2}}

% Constant commands
\newcommand{\avogadro}{6.022\times10^{23}}

\newcommand{\boltzmannConstK}{1.381\times10^{-23}\frac{\text{J}}{\text{K}}}
\newcommand{\boltzmannConstInEV}{8.617\times10^{-5}\frac{\text{eV}}{\text{K}}}
\newcommand{\planckConst}{6.626\times10^{-34}\text{J}\cdot\text{s}}

\newcommand{\atmP}{1.01325\times10^5\,\text{Pa}}

% Commands to easily write partial derivatives
\newcommand{\partDeriv}[2]{\frac{\partial #1}{\partial #2}}
\newcommand{\doublePartDeriv}[2]{\frac{\partial^2 #1}{\partial^2 #2}}


\title{Problem 7.11}
\author{Gabriel Decker}


\begin{document}



\maketitle
\begin{flushleft}


For this problem, we will determine the probability of a fermion having a particular state, varying the energy levels.
To do this, we'll use the fermi-dirac distribution.
The fermi-dirac distribution technically gives the average number of particles in a state, but since it lies between 0 and 1, it also acts as a probability.\\~\\

The distribution is the following.
\begin{equation}
    \overline{n}_\text{FD}=\frac{1}{e^{(\epsilon-\mu)/kT}+1}
\end{equation}
where $\epsilon$ is the energy of the state.\\~\\

We'll find the probabilities for the following values of $\epsilon$.
\begin{enumerate}
    \item[a.] $\epsilon=\mu-1\,\text{eV}$:\\
    \item[b.] $\epsilon=\mu-0.01\,\text{eV}$:\\
    \item[c.] $\epsilon=\mu$:\\
    \item[d.] $\epsilon=\mu+0.01\,\text{eV}$:\\
    \item[e.] $\epsilon=\mu+1\,\text{eV}$:\\
\end{enumerate}

Before plugging in numbers, let's make a general formula, representing $\epsilon$'s offset from $\mu$ as $\delta$.
$$\mathcal{P}(\delta)=\frac{1}{e^{(\epsilon-\mu)/kT}+1}$$
$$\implies\mathcal{P}(\delta)=\frac{1}{e^{((\mu+\delta)-\mu)/kT}+1}$$
$$\implies\mathcal{P}(\delta)=\frac{1}{e^{\delta/kT}+1}$$
We'll use room temp as $T=300\,\text{K}$ and $k=\boltzmannConstInEV$.\\~\\

\textbf{Part a}\\
$\delta=-1\,\text{eV}$. Plugging this in, we get the following probability.
$$\mathcal{P}(\delta)=\frac{1}{e^{\delta/kT}+1}$$
$$\implies\mathcal{P}(\delta)=\left(e^{(-1\,\text{eV})/((\boltzmannConstInEV)\cdot(300\,\text{K}))}+1\right)^{-1}$$
$$\implies\mathcal{P}(\delta)=1.000$$\\~\\

\textbf{Part b}\\
$\delta=-0.01\,\text{eV}$. Plugging this in, we get the following probability.
$$\mathcal{P}(\delta)=\frac{1}{e^{\delta/kT}+1}$$
$$\implies\mathcal{P}(\delta)=\left(e^{(-0.01\,\text{eV})/((\boltzmannConstInEV)\cdot(300\,\text{K}))}+1\right)^{-1}$$
$$\implies\mathcal{P}(\delta)=0.596$$\\~\\

\textbf{Part c}\\
$\delta=0\,\text{eV}$. Plugging this in, we get the following probability.
$$\mathcal{P}(\delta)=\frac{1}{e^{\delta/kT}+1}$$
$$\implies\mathcal{P}(\delta)=\left(e^{(0\,\text{eV})/((\boltzmannConstInEV)\cdot(300\,\text{K}))}+1\right)^{-1}$$
$$\implies\mathcal{P}(\delta)=\frac{1}{2}$$\\~\\

\textbf{Part d}\\
$\delta=0.01\,\text{eV}$. Plugging this in, we get the following probability.
$$\mathcal{P}(\delta)=\frac{1}{e^{\delta/kT}+1}$$
$$\implies\mathcal{P}(\delta)=\left(e^{(0.01\,\text{eV})/((\boltzmannConstInEV)\cdot(300\,\text{K}))}+1\right)^{-1}$$
$$\implies\mathcal{P}(\delta)=0.404$$\\~\\

\textbf{Part e}\\
$\delta=1\,\text{eV}$. Plugging this in, we get the following probability.
$$\mathcal{P}(\delta)=\frac{1}{e^{\delta/kT}+1}$$
$$\implies\mathcal{P}(\delta)=\left(e^{(1\,\text{eV})/((\boltzmannConstInEV)\cdot(300\,\text{K}))}+1\right)^{-1}$$
$$\implies\mathcal{P}(\delta)=0.000$$\\~\\




\end{flushleft}

\end{document}
